\section{Metodologia}

\subsection{Fonte e descrição dos dados}

Os dados utilizados neste trabalho são provenientes do portal eSMR (Electronic Self-Monitoring Report) do California Water Boards, referentes à estação de tratamento Los Angeles--Glendale Water Reclamation Plant (LAGWRP). Cada linha do relatório corresponde a uma medição ou a um valor derivado (carga mássica, média diária/semanal/mensal) de um parâmetro de qualidade da água em um ponto de monitoramento, com colunas de local (\texttt{Location}), parâmetro (\texttt{Parameter}), método (\texttt{Analytical Method}/\texttt{Calculated Method}), qualificador (\texttt{Qual}), resultado (\texttt{Result}), unidade (\texttt{Units}) e limites analíticos (\texttt{MDL}, \texttt{ML}, \texttt{RL}), entre outras.

Foram utilizados quatro parâmetros: Sólidos Dissolvidos Totais (TDS), Cloreto, Amônia Total (como N) e Demanda Bioquímica de Oxigênio (BOD, 5 dias a 20\textdegree{}C). Os dois pontos de código \texttt{EFF-001} (até 2012) e \texttt{EFF-001A} (a partir de 2012) correspondem ao mesmo ponto físico de monitoramento -- o efluente final da estação -- e foram unificados em uma única série temporal.

A série mensal canônica foi construída filtrando, para cada parâmetro: local em \{\texttt{EFF-001}, \texttt{EFF-001A}\}; \texttt{Calculated Method} igual a "Monthly Average (Mean)" (a média mensal já pré-calculada e reportada pela concessionária); e \texttt{Units} igual a "mg/L" (concentração, unidade relevante para os objetivos deste trabalho -- a carga mássica em lb/day não foi utilizada na análise de tendência/correlação). O período resultante cobre de fevereiro de 2011 a março de 2026 (182 meses), com um único mês (julho de 2011) sem valor de BOD reportado pela concessionária.

\subsection{Tratamento dos valores não detectados (ND) de BOD}

Das 182 médias mensais de BOD no efluente, 118 (65\%) são reportadas pela concessionária como não detectadas (\texttt{Qual = ND}, resultado em branco), com limite de detecção do método (MDL) constante em 3,0~mg/L ao longo de todo o período. As 63 médias mensais com valor detectado concentram-se muito próximas desse mesmo piso (média de 3,08~mg/L, mediana 3,0~mg/L, máximo 4,0~mg/L), indicando que a concentração de BOD no efluente permanece rente ao limite de detecção do método durante os 15 anos observados.

Uma verificação sistemática dos dados brutos diários de BOD (5.422 amostras individuais, das quais 3.205 são ND) mostrou que, nos casos em que existe uma linha de carga mássica diária (\texttt{Daily Discharge}, lb/day) correspondente, essa carga é sempre igual a zero quando a concentração é ND (316 de 316 casos) e sempre maior que zero quando a concentração é detectada (2.217 de 2.217 casos) -- uma convenção interna consistente, embora disponível para apenas cerca de 10\% dos dias ND. Verificou-se ainda que não há correlação relevante entre a taxa de ND mensal do BOD e o valor de TDS do mesmo mês (correlação ponto-bisserial $r = 0{,}03$, $p = 0{,}66$), o que reduz o risco de a escolha de tratamento induzir artificialmente uma correlação TDS--BOD.

Diante da ambiguidade metodológica -- convenção interna sugerindo zero, prática científica padrão para dados censurados à esquerda recomendando MDL/2 -- optou-se por não escolher um único tratamento a priori. Foi avaliada ainda uma terceira opção, informada por Antweiler e Taylor \cite{antweiler2008evaluation}: ao reavaliar as principais técnicas para dados censurados à esquerda (substituição, máxima verossimilhança, Kaplan-Meier e ROS) usando conjuntos de dados reais em que o valor verdadeiro dos dados censurados era conhecido, os autores concluem que, para conjuntos com menos de 70\% de censura (o caso do nosso BOD, com 65\%), Kaplan-Meier é a melhor técnica geral para \textbf{estatísticas-resumo} (média, mediana, quartis), com ROS e os métodos de substituição por MDL/2 ou por valor aleatório entre 0 e o MDL como alternativas adequadas. Como Kaplan-Meier estima a distribuição da amostra como um todo (não um valor individual por observação), ele não se presta a gerar, por si só, um valor substituto para cada uma das 118 observações ND necessárias para compor a série mensal de BOD deste trabalho -- por isso a técnica efetivamente aplicada aqui é o ROS (\textit{Regression on Order Statistics}, a "alternativa adequada" indicada pelos mesmos autores): ajusta uma reta de regressão entre os valores detectados e seus quantis normais esperados (posição de plotagem de Hazen, assumindo que as observações censuradas ocupam os postos mais baixos) e extrapola o valor esperado da fração censurada a partir da forma real da distribuição observada, em vez de uma convenção arbitrária. O ajuste tem $r=0{,}677$ ($p=1{,}1\times10^{-9}$) e produz um valor substituto estimado de 2,517~mg/L -- entre o MDL/2 (1,5) e o MDL cheio (3,0), sugerindo que a convenção MDL/2 pode estar subestimando sistematicamente o BOD real dos meses ND. Por construção, o ROS também produz um único valor substituto para todas as 118 observações ND (indistinguíveis entre si, todas reportadas apenas como "< 3,0"): não recupera variação mês a mês real dentro dos meses censurados, mesma limitação estrutural das opções A e B, apenas com um nível melhor embasado nos dados observados. Não foi encontrado padrão temporal monotônico na taxa de ND (oscila entre 25\% e 100\% ao ano, sem tendência consistente de alta ou queda ao longo dos 15 anos).

Foram construídos três datasets mensais canônicos, idênticos em TDS, Cloreto e Amônia, diferindo apenas no tratamento do BOD ND:

\begin{itemize}
    \item \textbf{Dataset A} (\texttt{dataset\_canonico\_bod\_mdl2.csv}): valores ND substituídos por MDL/2 (1,5~mg/L);
    \item \textbf{Dataset B} (\texttt{dataset\_canonico\_bod\_zero.csv}): valores ND substituídos por zero;
    \item \textbf{Dataset F} (\texttt{dataset\_canonico\_bod\_ros.csv}): valores ND substituídos pela estimativa ROS/Helsel (2,517~mg/L).
\end{itemize}

Como TDS, Cloreto e Amônia independem do tratamento de ND do BOD, apenas a análise de correlação TDS--BOD (Seção de Resultados) é sensível a essa escolha e roda nos três datasets em paralelo -- os demais métodos da bateria (univariados em TDS) produzem o mesmo resultado nos três, e por isso rodam uma única vez.

\subsection{Bateria de métodos}

Serão aplicados, de forma independente e em paralelo, os seguintes métodos sobre a série mensal canônica: (a) estatística clássica de tendência -- Mann-Kendall com inclinação de Sen, Theil-Sen e regressão OLS; (b) séries temporais clássicas -- decomposição STL e ARIMA/SARIMA; (c) modelos baseados em árvore -- Random Forest e XGBoost/LightGBM; (d) Prophet e regressão bayesiana, que expõem explicitamente a incerteza crescente da previsão em horizontes longos; (e) quatro \textit{baselines} obrigatórios -- naive, naive sazonal, ETS/Holt-Winters e Theta -- contra os quais todo método mais sofisticado precisa ser julgado antes de ser considerado um resultado válido. Cada método produzirá previsões pontuais e intervalos de confiança de 90\% para +10, +15 e +20 anos a partir do último dado observado, além de métricas de desempenho em holdout (RMSE, MAE, R\textsuperscript{2}) e a taxa de tendência estimada (mg/L/ano) com respectivo p-valor. Esta seção será atualizada com os resultados reais assim que cada método for executado (ver Seção de Resultados).

\subsection{Análise de estrutura da série}

Antes de qualquer ajuste de modelo, a série de TDS foi submetida a um diagnóstico estrutural (\texttt{script\_07\_analise\_estrutura\_serie.py}), para não assumir propriedades (sazonalidade, estacionariedade, homogeneidade) sem testá-las: (i) decomposição STL e a estatística de força de sazonalidade de Wang et al. ($F_s = \max(0, 1 - \mathrm{Var}(\mathrm{residuo})/\mathrm{Var}(\mathrm{sazonal}+\mathrm{residuo}))$), comparada ao limiar de referência $F_s > 0{,}64$ da literatura de \textit{forecasting}; (ii) testes de estacionariedade ADF (hipótese nula: raiz unitária) e KPSS (hipótese nula: estacionariedade) -- hipóteses opostas por desenho, reportadas mesmo quando aparentemente contraditórias; (iii) testes de quebra estrutural: Chow (com \textit{breakpoint} fixo na data aproximada da troca de código \texttt{EFF-001}$\to$\texttt{EFF-001A}, $\sim$2012), Pettitt (não-paramétrico, sem \textit{breakpoint} pré-definido, implementado diretamente pois não está disponível em \texttt{pymannkendall}) e CUSUM dos resíduos de uma tendência linear simples (\texttt{statsmodels.stats.diagnostic.breaks\_cusumolsresid}). O resultado da força de sazonalidade também condiciona se o termo sazonal do ETS/Holt-Winters e a etapa de dessazonalização do Theta (bateria de \textit{baselines}) são ativados.

\subsection{Framework de validação temporal honesta}

Para além do holdout único dos últimos 24 meses já usado desde a primeira versão da bateria, cada método (\texttt{script\_01} a \texttt{script\_05} e \texttt{script\_08}) passa a reportar, via módulo compartilhado \texttt{validacao\_utils.py}: (i) MASE (\textit{Mean Absolute Scaled Error}, escalado pelo erro do naive sazonal $m=12$ no próprio treino) e sMAPE, métricas de erro escala-independente que permitem comparar métodos entre si e contra os \textit{baselines}; (ii) validação cruzada de séries temporais com janela expansiva (5 \textit{folds}, nunca \textit{split} aleatório, nunca vazamento temporal -- o treino de cada \textit{fold} é sempre estritamente anterior ao teste); (iii) \textit{backtesting} com origem móvel: como nenhum ponto do histórico de $\sim$15 anos tem de fato 10--20 anos de futuro observado, o \textit{backtest} mede o erro real em horizontes de +3 e +5 anos a partir de várias origens de treino ao longo da série -- o único proxy honesto de desempenho de extrapolação disponível para os horizontes realmente pedidos pelo projeto. Por restrição de custo computacional, os hiperparâmetros de métodos com busca em grade cara (SARIMA, Random Forest, XGBoost, LightGBM) são fixados nos valores já escolhidos pela busca original em vez de repetidos a cada um dos $\sim$25 \textit{folds}/origens de validação -- decisão de escopo documentada nos próprios scripts, não escondida. Pelo mesmo motivo, a validação de \textit{fold} da regressão bayesiana usa o estimador de mínimos quadrados equivalente (a média da posterior de um modelo linear-gaussiano com \textit{priors} fracos converge para a OLS) em vez de reamostrar NUTS 25 vezes; a previsão final (+10/+15/+20 anos) reportada em \texttt{resultados\_comparacao.csv} continua usando NUTS completo, com incerteza posterior real. Todos os métodos desta e das subseções seguintes (novos métodos) são submetidos ao mesmo framework -- MASE holdout, CV expansiva e \textit{backtest} de origem móvel são reportados uniformemente para as 21 linhas de \texttt{resultados\_comparacao.csv}.

\subsection{Métodos adicionais}

Além da bateria original e dos \textit{baselines}, cinco métodos adicionais foram testados, escolhidos livremente a partir do material de apoio de referências para atacar limitações já observadas empiricamente nos métodos anteriores:

\textbf{SVR e Gaussian Process} (\texttt{script\_09\_svr\_gp.py}): SVR (\textit{kernel} RBF, hiperparâmetros por \texttt{GridSearchCV}+\texttt{TimeSeriesSplit}, mesmas \textit{features} de tempo/sazonalidade/\textit{lags} de Random Forest, IC90 por \textit{bootstrap} de resíduos -- 200 réplicas da previsão recursiva injetando um resíduo do treino a cada passo). Gaussian Process com \textit{kernel} composto (RBF para tendência de longo prazo + \texttt{ExpSineSquared} periódico de 12 meses + \texttt{WhiteKernel} para ruído), ajustado diretamente sobre o tempo, IC90 analítico via desvio-padrão posterior.

\textbf{Tendência + árvore no resíduo} (\texttt{script\_10\_detrend\_arvore.py}): ataca diretamente a limitação de extrapolação de árvores já documentada (Random Forest, XGBoost e LightGBM saturam ou oscilam na previsão de longo prazo). Ajusta uma tendência OLS na série completa, treina Random Forest e XGBoost para prever o \textbf{resíduo} (série sem tendência, dentro do range de treino mesmo longe no futuro) e soma a tendência extrapolada de volta à previsão do resíduo. A incerteza da própria tendência OLS não é propagada ao IC90 final (que vem só da dispersão do resíduo previsto) -- subestima a incerteza total, limitação registrada.

\textbf{SARIMAX multivariado com Cloreto} (\texttt{script\_11\_multivariado\_cloreto.py}): adiciona Cloreto como regressor exógeno ao SARIMA (ordem reaproveitada da busca por AIC de \texttt{script\_02}, por custo computacional), testando se o preditor amplamente usado na literatura (material de apoio, Tema 2.4) melhora a previsão de TDS neste dataset. Holdout/CV/\textit{backtest} usam o Cloreto historicamente observado no período de teste (avaliação honesta de "se o Cloreto futuro fosse conhecido"); a previsão real de +10/+15/+20 anos exige extrapolar o próprio Cloreto (via Theil-Sen), cuja incerteza não é propagada ao IC90 do TDS -- limitação explícita.

\textbf{Híbrido SARIMA+Prophet} (\texttt{script\_12\_hibrido\_arima\_prophet.py}): \textit{ensemble} pela média aritmética das previsões pontuais de SARIMA e Prophet (mecanismos de extrapolação distintos, já mostrados divergentes em direções opostas nos resultados anteriores), com IC90 pela envoltória dos dois intervalos individuais.

\textbf{LSTM leve} (\texttt{script\_13\_deep\_learning.py}, PyTorch CPU): 1 camada, 32 unidades, janela de 12 meses de \textit{lookback} -- arquitetura deliberadamente pequena, já que $\sim$180 pontos mensais não sustentam uma rede maior sem sobreajuste severo. Holdout avaliado em modo \textit{one-step} (mesma convenção de RF/XGBoost/SVR); previsão recursiva com IC90 por \textit{bootstrap} de resíduos para os horizontes longos e para CV/\textit{backtest}. Testado explicitamente sob o mesmo framework de validação honesta -- o material de apoio (Tema 6.3) documenta um precedente (XGBoost venceu LSTM em vazão mensal) tratado aqui como expectativa a testar, não verdade assumida; o resultado real decide se o método é reportado como finalista.

\subsection{Reconstrução da vazão do efluente}

O dataset traz o mesmo parâmetro em duas unidades -- \texttt{mg/L} (concentração) e \texttt{lb/day} (carga mássica) -- relacionadas pela identidade padrão do setor $\mathrm{lb/day} = \mathrm{mg/L} \times \mathrm{vazao(MGD)} \times 8{,}34$, o que permite reconstruir a vazão do efluente da LAGWRP, não disponível explicitamente no dataset (\texttt{script\_16\_reconstrucao\_vazao.py}). A validação usou duas frentes: plausibilidade contra a capacidade nominal de 20~MGD da planta (brochure institucional), e consistência cruzada entre TDS, Cloreto, Amônia e BOD -- medidos na mesma amostra física, a vazão derivada independentemente de cada um deve coincidir se a identidade for real. A vazão derivada de TDS correlaciona 0,997 e 0,998 com a de Cloreto e Amônia respectivamente (0,87 com BOD, mais ruidosa mas ainda forte), com média de 9,49~MGD ($\sim$47\% da capacidade nominal) e nenhum mês fora de uma faixa plausível de 2-40~MGD -- a identidade se sustenta nos dados reais. Isso habilita os métodos WRTDS, balanço de massa e simulação de cenários como extensão futura da bateria, condicionada à aprovação do usuário (Seção de Resultados/Conclusão).

\subsection{Matriz de sensibilidade no tratamento de dados}

Antes de fixar qualquer tratamento de dados como padrão, 9 decisões de pré-processamento com risco conhecido de viés (\texttt{script\_17\_matriz\_sensibilidade.py}) foram testadas: tratamento de ND do BOD (5 variantes -- MDL/2, zero, Kaplan-Meier, ROS/Helsel, máxima verossimilhança com distribuição log-normal censurada), quebra de código \texttt{EFF-001}$\to$\texttt{EFF-001A}, mudanças de MDL/método analítico ao longo do tempo, reagregação a partir de amostras brutas vs.\ uso da média mensal pré-calculada, agregação anual (média/mediana $\times$ ano civil/hidrológico), tratamento de \textit{outliers}, meses faltantes, correção de autocorrelação no teste de Mann-Kendall (\textit{pre-whitening}, \textit{trend-free pre-whitening}, correção de variância de Hamed-Rao, Seasonal Kendall -- todas via \texttt{pymannkendall}), e escala bruta vs.\ logarítmica. Os resultados (\texttt{matriz\_sensibilidade\_resultados.csv}) mostraram que a tendência de alta de TDS é sensível a três dessas escolhas -- decisão de tratamento padrão pendente de aprovação do usuário, e por isso a tabela completa ainda não foi transcrita para a Seção de Resultados (evita reescrevê-la duas vezes).
